%%%%%%%%%%%%%%%%%%%%%%
%%%%%%%%%%%%%%%%%%%%%%

\begin{figure*}
\begin{centering}
	\begin{subfigure}[b]{0.45\textwidth}
                \centering
                \includegraphics[width=\textwidth]{Figures/oc/new_results/human_mouse_bpo_2}
                \caption{Human/Mouse}                                
                \label{fig:ocHMB2}
        \end{subfigure}%
	\quad
	\begin{subfigure}[b]{0.45\textwidth}
                \centering
                \includegraphics[width=\textwidth]{Figures/oc/new_results/human_rat_bpo_2}
                \caption{Human/Rat}                                
                \label{fig:ocHRB2}
        \end{subfigure}%

	\begin{subfigure}[b]{0.45\textwidth}
                \centering
                \includegraphics[width=\textwidth]{Figures/oc/new_results/mouse_rat_bpo_2}
                \caption{Mouse/Rat}                                
                \label{fig:ocMRB2}
        \end{subfigure}%
	\quad
	\begin{subfigure}[b]{0.45\textwidth}
                \centering
                \includegraphics[width=\textwidth]{Figures/oc/new_results/worm_fly_bpo_2}
                \caption{Worm/Fly}                                
                \label{fig:ocWFB2}
        \end{subfigure}%



\par\end{centering}
\caption[Comparison of orthologs and paralogs in four pairs of organisms for Biological Process annotations]{Comparison of orthologs (red) and paralog (blue) functional similarity across a range of global sequence identity values in four pairs of organisms when considering Biological Process annotations. Standard error bars are shown. }
\label{fig:oc1_bpo}
\end{figure*}

%%%%%%%%%%%%%%%%%%%%%%
%%%%%%%%%%%%%%%%%%%%%%